\maketitle
\makesignature

\ifproject
\begin{abstractTH}
\enskip ปัจจุบันผู้ใช้บริการฟิตเนส โดยเฉพาะผู้เริ่มต้น มักประสบปัญหาขาดความรู้ในการออกกำลังกาย ขาดแรงจูงใจ และไม่มีระบบติดตามความก้าวหน้า ขณะเดียวกันผู้ประกอบการฟิตเนสยังขาดเครื่องมือดิจิทัลที่ช่วยบริหารจัดการเครื่องเล่นและสร้างการมีส่วนร่วมกับสมาชิก โครงงานนี้มีวัตถุประสงค์เพื่อสำรวจและออกแบบระบบ GymMate ซึ่งเป็นแพลตฟอร์มในรูปแบบ Web/Mobile Application ที่ผสานแนวคิด Gamification เข้ากับ Workout Recommendation และ Program Planner ระบบที่ออกแบบมีฟังก์ชันสำคัญ ได้แก่ โปรแกรมการฝึก (ทั้งแบบสำเร็จรูปและแบบปรับแต่งได้), การแสดงแผนผังยิม (Floorplan) พร้อมปุ่ม Flex Substitute เพื่อแนะนำเครื่องหรือท่าทางทดแทนเมื่ออุปกรณ์ไม่ว่าง, ระบบบันทึกและวิเคราะห์ความก้าวหน้า, การสะสมคะแนนและจัดอันดับเพื่อสร้างแรงจูงใจ รวมถึงแดชบอร์ดสำหรับผู้ดูแลในการจัดการข้อมูลเครื่องเล่นและพื้นที่ยิม ผลลัพธ์ของโครงงานนี้คือชุดคุณลักษณะและต้นแบบ UX/UI ที่ผ่านการเก็บและวิเคราะห์ความต้องการจากอาจารย์ที่ปรึกษา เจ้าของยิม และผู้จัดการฟิตเนส โดยอ้างอิงมาตรฐาน ACSM Guidelines และ Physical Activity Guidelines เพื่อนำไปสู่การพัฒนาระบบต้นแบบที่สามารถสนับสนุนการออกกำลังกายอย่างมีประสิทธิภาพและยั่งยืนได้ในอนาคต
\end{abstractTH}



\begin{abstract}
Fitness users, especially beginners, often face challenges such as a lack of exercise knowledge, insufficient motivation, and the absence of effective progress-tracking tools. At the same time, gym owners lack digital solutions to manage equipment efficiently and engage members. This project aims to survey and design GymMate, a Web/Mobile application platform that integrates Gamification with Workout Recommendation and Program Planning. The proposed system includes key features such as preset and customizable workout programs, an interactive Gym Floorplan with a Flex Substitute function to suggest alternative machines or exercises when equipment is occupied, progress tracking and analytics, gamification elements (points, badges, streaks, and leaderboards), and an admin dashboard for managing gym equipment and information. The outcome of this project is a refined set of features and UX/UI prototypes, derived through iterative requirement gathering with advisors, gym owners, and fitness managers. The design also references the ACSM Guidelines and Physical Activity Guidelines, laying the foundation for a future prototype that supports effective and sustainable fitness training.
\end{abstract}

\iffalse
\begin{dedication}
This document is dedicated to all Chiang Mai University students.

Dedication page is optional.
\end{dedication}
\fi % \iffalse

\begin{acknowledgments}
โครงงาน \textbf{GymMate : เพื่อนรู้ใจเรื่องฟิตเนส (GymMate : Your Fitness Buddy)} ไม่อาจสำเร็จลุล่วงได้ หากปราศจากแรงสนับสนุน คำแนะนำ และความช่วยเหลือจากหลายฝ่าย คณะผู้จัดทำขอถือโอกาสนี้แสดงความขอบคุณอย่างจริงใจต่อทุกท่านที่มีส่วนร่วมในการผลักดันให้โครงงานนี้สำเร็จลุล่วง

คณะผู้จัดทำขอขอบคุณ \textbf{ผู้ช่วยศาสตราจารย์ ดร.\ นวดนย์ คุณเลิศกิจ} อาจารย์ที่ปรึกษาโครงงาน ที่ได้ให้คำแนะนำทั้งด้านวิชาการและแนวทางการพัฒนาอย่างต่อเนื่องตลอดระยะเวลาของโครงงาน นอกจากนี้ขอขอบคุณ \textbf{ผู้ช่วยศาสตราจารย์ โดม โพธิกานนท์} และ \textbf{ผู้ช่วยศาสตราจารย์ ดร.\ ลัชนา ระมิงค์วงศ์} อาจารย์กรรมการสอบโครงงาน ที่ได้สละเวลาตรวจสอบและให้ข้อเสนอแนะอันมีค่า

คณะผู้จัดทำขอขอบคุณ \textbf{ภาควิชาวิศวกรรมคอมพิวเตอร์ คณะวิศวกรรมศาสตร์ มหาวิทยาลัยเชียงใหม่} ที่ให้การสนับสนุนด้านทรัพยากรและองค์ความรู้ รวมถึงผู้เข้าร่วมทดสอบการใช้งานระบบ (Usability Testing) ทุกท่าน ที่ให้ข้อเสนอแนะอันเป็นประโยชน์ต่อการพัฒนาระบบ สุดท้ายขอขอบคุณครอบครัวทุกท่านที่ให้การสนับสนุนและเป็นกำลังใจตลอดมา

\acksign{2026}{3}{24}
\end{acknowledgments}
\fi % \ifproject

\contentspage

\ifproject
\figurelistpage

\tablelistpage
\fi % \ifproject

% \abbrlist % this page is optional

% \symlist % this page is optional

% \preface % this section is optional
